
\section{امتیازی: نا‌بودن کران دودلی}
(امتیازی) در این تمرین قصد داریم نشان دهیم که نامساوی دودلی در یک مسئله‌ی خاص،
\lr{tight}
نیست. فرض کنید 
$e_1, e_2, \dots, e_n$
پایه‌ی استاندارد (کانونی) فضای
$\mathbb{R}^n$
باشد. مجموعه‌ی $T$ را به شکل زیر تعریف می‌کنیم:

\begin{equation*}
T = \Big\{\frac{e_k}{\sqrt{1 + \log(k)}}, \; k=1,\dots,n\Big\}
\end{equation*}

نشان دهید:

\begin{itemize}
	\item 
	نامساوی زیر برای عرض گوسی مجموعه‌ی $T$ محدود است.
	\item 
	نشان دهید انتگرال دودلی واگراست، یعنی
	$$\int_{0}^{\infty}\sqrt{\log(\mathcal{N}(T, d, \epsilon))} d\epsilon \to \infty.$$
\end{itemize}
